\chapter{The Decoherent Hegemon (USA): Iper-estensione senza
Coerenza}\label{legemone-decoerente-usa-hyper-extension-senza-coerenza}

\begin{quote}
\textbf{Core argument:} the United States del secondo mandato Trump
esibiscono il pattern formal dell'\emph{egemone decoerente}: massima
material capacity \(s_{ii}\) residua, coupling sistemico interno
\(K > K_c\), ma order parameter \(\Phi \to 0\) per effetto combinato
di hyper-extension esecutiva, affective polarization, declino di
\(s_{iii}\) e azione di sorgenti esogene di disordine \(\zeta(t)\) ---
fra cui il Soliton Musk del Ch.~6. Il sistema è formalmente in regime
sub-critico despite il superamento della soglia di coupling.
\end{quote}



\section{Structural Definition dell'Egemone
Decoerente}\label{definition-structural-dellegemone-decoerente}

\Hypoth{} \textbf{Definizione (Decoherent Hegemon DEQ$^2$).} Un attore
\(e\) è detto \emph{Decoherent Hegemon} se soddisfa simultaneamente:

\begin{enumerate}
\def\labelenumi{\arabic{enumi}.}
\tightlist
\item
  \(s_{ii,e} \geq \langle s_{ii} \rangle_{\text{sis}} + \sigma_{s_{ii}}\)
  (material capacity residua sopra la media sistemica per almeno una
  deviazione standard);
\item
  \(K_e > K_c\) (coupling di Kuramoto interno superiore alla soglia
  critica --- esistono \emph{meccanismi formali} di coordinazione:
  partiti, federalismo, mercato, media);
\item
  \(\Phi_e \to 0\) malgrado (1) e (2) --- il order parameter collassa
  per \emph{azione combinata} di rumore endogeno (polarizzazione
  affettiva), perdita di legitimacy (\(s_{iii} \downarrow\)) e sorgenti
  esogene \(\zeta(t)\);
\item
  \(A_e < 0\) a livello sistemico (la risposta a \(\sigma\) è concava:
  la volatilità non viene assorbita ma amplificata dall'iper-reattività
  esecutiva);
\item
  \(\Omega_e \uparrow\) (energia istituzionale dissipata in conflitto
  interno cresce più rapidamente dell'energia utile prodotta).
\end{enumerate}

\Proved{} \textbf{Confollowsnza.} \(T_s^{(2),\text{sys}} \to 0\) per
\(\Phi \to 0\) al numeratore e \(\Omega/\Phi \to \infty\) al
denominatore. La phase transition è \textbf{structural}, non
congiunturale.



\section{\texorpdfstring{Misura Empirica della Decoerenza:
\(\Phi^{\text{USA}}\) tra 2024 e
2026}{Misura Empirica della Decoerenza: \textbackslash Phi\^{}\{\textbackslash text\{USA\}\} tra 2024 e 2026}}\label{misura-empirica-della-decoherence-phitextusa-tra-2024-e-2026}

\Proved{} \textbf{Indicatori convergenti del collapse di \(\Phi\).}

{\def\LTcaptype{none} % do not increment counter
{\small\begin{longtable}[]{@{}
  >{\raggedright\arraybackslash}p{(\linewidth - 6\tabcolsep) * \real{0.2500}}
  >{\raggedright\arraybackslash}p{(\linewidth - 6\tabcolsep) * \real{0.2500}}
  >{\raggedright\arraybackslash}p{(\linewidth - 6\tabcolsep) * \real{0.2500}}
  >{\raggedright\arraybackslash}p{(\linewidth - 6\tabcolsep) * \real{0.2500}}@{}}
\toprule\noalign{}
\begin{minipage}[b]{\linewidth}\raggedright
Indicatore
\end{minipage} & \begin{minipage}[b]{\linewidth}\raggedright
Valore 2024
\end{minipage} & \begin{minipage}[b]{\linewidth}\raggedright
Valore 2026
\end{minipage} & \begin{minipage}[b]{\linewidth}\raggedright
Variazione
\end{minipage} \\
\midrule\noalign{}
\endhead
\bottomrule\noalign{}
\endlastfoot
V-Dem Liberal Democracy Index USA & \(\approx 0{,}73\) & \(0{,}57\)
(minimo dal 1965) & \(-22\%\) \\
Classificazione V-Dem & Liberal Democracy & Electoral Democracy
(declassato) & salto categoriale \\
Constraint legislativi/giudiziari sull'esecutivo & sopra media G7 &
minimo da 100+ anni & crollo storico \\
Trust nel governo federale (Pew) & Dem \(\approx 18\%\) & Dem \(9\%\)
(minimo storico), Rep \(26\%\) & bipartisan disaffection \\
Quota cittadini che vede l'altro partito come \emph{minaccia} &
\(\approx 35\%\) & \(> 40\%\) & \(+5\) pt in due anni \\
Common ground percepito tra partiti (sei aree, Pew) & baseline 2023 &
\(-12\) pt medi & dissoluzione del centro \\
\end{longtable}}
}

\Hypoth{} \textbf{Interpretazione formal.} L'indice V-Dem è proxy
diretta di \(s_{iii}\) (ethical-discursive legitimacy) ed indiretta di
\(\Phi\) (la \emph{liberal democracy} esige synchronization interna fra
esecutivo, legislativo, giudiziario, civil society). Il collapse
simultaneo di V-Dem e di Pew Trust è la firma empirica di
\(\Phi \to 0\): i sotto-attori (i cittadini, raggruppati per partito)
hanno fasi \(\theta_j\) sempre più distanti, e l'aggregato
\(|N^{-1}\sum_j e^{i\theta_j}|\) tende a zero per cancellazione di
interferenza.



\section{\texorpdfstring{Iper-estensione Esecutiva:
\(A_{\text{sys}} < 0\)}{Iper-estensione Esecutiva: A\_\{\textbackslash text\{sys\}\} \textless{} 0}}\label{hyper-extension-esecutiva-a_textsys-0}

\Hypoth{} \textbf{Tesi.} Il pattern dell'amministrazione Trump 2.0 nel
primo anno (gennaio 2025 --- gennaio 2026) configure ciò che la
letteratura tecnica chiama \emph{escalation as mode of governance}: 253
Executive Orders, 59 memoranda e 135 proclamations al 31 marzo 2026 ---
un volume normativo unilaterale senza precedenti nella storia
repubblicana.

\Proved{} \textbf{Confollowsnza formal per \(A_{\text{sys}}\).}
L'antifragility è definita come
\(\partial^2 \Pi / \partial \sigma^2 > 0\). Un sistema che reagisce ad
ogni shock con concentrazione esecutiva \emph{riduce strutturalmente} lo
spazio dei meccanismi di assorbimento distribuito (federalismo, judicial
review, civil service indipending). La risposta diventa concava: ogni
nuovo shock richiede un'escalation maggiore del precedente per produrre
lo stesso effetto. Formalmente
\(\partial^2 \Pi / \partial \sigma^2 < 0\).

\Conj{} \textbf{Caveat metodologico.} \(A_e\) individuale del Presidente
(alta opzionalità, alta tolleranza al rischio) può restare positivo ---
è il livello \emph{sistemico} che diventa fragile. Si replica a scala
macro la stessa asimmetria del Ch.~6: opzionalità individuale \(\neq\)
antifragility collettiva.

\Proved{} \textbf{Verifica indiretta dei vincoli giudiziari.}

{\def\LTcaptype{none} % do not increment counter
{\small\begin{longtable}[]{@{}
  >{\raggedright\arraybackslash}p{(\linewidth - 4\tabcolsep) * \real{0.3333}}
  >{\raggedright\arraybackslash}p{(\linewidth - 4\tabcolsep) * \real{0.3333}}
  >{\raggedright\arraybackslash}p{(\linewidth - 4\tabcolsep) * \real{0.3333}}@{}}
\toprule\noalign{}
\begin{minipage}[b]{\linewidth}\raggedright
Evento
\end{minipage} & \begin{minipage}[b]{\linewidth}\raggedright
Data
\end{minipage} & \begin{minipage}[b]{\linewidth}\raggedright
Effetto su \(A_{\text{sys}}\)
\end{minipage} \\
\midrule\noalign{}
\endhead
\bottomrule\noalign{}
\endlastfoot
\emph{Trump v. CASA} --- limita injunctions nazionali & 27 giugno 2025 &
\(A_{\text{sys}} \downarrow\) (rimuove un assorbitore distribuito) \\
Sentenza tariffe IEEPA, 6-3 contro & 20 febbraio 2026 &
\(A_{\text{sys}} \uparrow\) marginale (vincolo attivato) \\
Risposta esecutiva: tariffa 10\% globale ex Section 122 & febbraio 2026
& escalation immediata --- bypass del vincolo \\
\end{longtable}}
}

\Conj{} \textbf{Predizione operativa.} Il pattern \emph{vincolo \(\to\)
escalation per via alternativa} è la firma DEQ$^2$ dell'hyper-extension: la
corte impone un vincolo, l'esecutivo riformula l'azione su base
statutaria diversa entro settimane. La frizione istituzionale si riduce
ad attrito puramente nominale.



\section{\texorpdfstring{Crollo di \(s_{iii}\): Soft Power
2026}{Crollo di s\_\{iii\}: Soft Power 2026}}\label{crollo-di-s_iii-soft-power-2026}

\Proved{} \textbf{Dato Brand Finance Global Soft Power Index 2026.} Gli
United States registrano \textbf{il calo più ripido tra tutti i 193 Stati
misurati}: \(-4{,}6\) punti su scala 100, da \(\approx 79{,}5\) a
\(74{,}9\). La China sale a \(73{,}5\), riducendo a meno di \(1{,}5\)
punti il vantaggio USA.

{\def\LTcaptype{none} % do not increment counter
{\small\begin{longtable}[]{@{}
  >{\raggedright\arraybackslash}p{(\linewidth - 4\tabcolsep) * \real{0.3333}}
  >{\raggedright\arraybackslash}p{(\linewidth - 4\tabcolsep) * \real{0.3333}}
  >{\raggedright\arraybackslash}p{(\linewidth - 4\tabcolsep) * \real{0.3333}}@{}}
\toprule\noalign{}
\begin{minipage}[b]{\linewidth}\raggedright
Sotto-indice
\end{minipage} & \begin{minipage}[b]{\linewidth}\raggedright
Variazione 2025-2026
\end{minipage} & \begin{minipage}[b]{\linewidth}\raggedright
Lettura DEQ$^2$
\end{minipage} \\
\midrule\noalign{}
\endhead
\bottomrule\noalign{}
\endlastfoot
Reputation & \(-11\) pt (crolla al 26$^\circ$) & proiezione di \(s_{iii}\)
esterna \\
Friendliness & \(-32\) & rifiuto del registro emancipativo \\
Generosity & \(-68\) & crollo della \emph{ethical quality} \(Q_0\)
macro \\
Ease of doing business & \(-21\) & crollo \(A\) percepito da terzi \\
Climate action support & \(-16\) & \(Q_0\) sull'asse \(z\) ---
dialectical complexity negata \\
Human rights & \(-10\) & crollo coordinata \(x\) --- formal coherence \\
Ethical standards & \(-4\) & erosione del baseline \(Q_0\) \\
\end{longtable}}
}

\Hypoth{} \textbf{Modellazione formal.} \(s_{iii}^{\text{USA}}\) è la
\emph{componente esportabile} del prodotto \(Q_0 \cdot M_e\) degli Stati
Uniti come emittente macro. La sequenza \(-32, -68, -21, -16\) non è
rumore: è la firma del passaggio dell'emittente USA
dall'\textbf{Equilibrio Dialogico} \((7, 6, 8)\) all'\textbf{Equilibrio
Egemonico} \((8, 9, 2)\) della Topografia App.~A --- lo stesso
equilibrio separante in cui si colloca il Soliton (Ch.~6.2). Lo Stato
egemone \emph{si solitonizza} discorsivamente, e perde thus la presa sui
Riceventi che ne avevano riconosciuto la legitimacy sotto il regime
precedente.

\Proved{} \textbf{Verifica empirica della tesi.} Il crollo di
\emph{Friendliness} (\(-32\)) e \emph{Generosity} (\(-68\)) sono i due
assi che misurano direttamente la dimensione \(z\) (complessità
dialettica, riconoscimento dell'altro). Il loro collapse simultaneo a
fronte della tenuta dell'asse \(y\) (capacità tecnica residua,
\(s_{ii}\) ancora alto) replica la firma \((8, 9, 2)\) del profilo
Egemonico-Solitonico.



\section{\texorpdfstring{La Frattura Tech-Policy come Realizzazione di
\(\zeta_{\text{Sol}}(t)\)}{La Frattura Tech-Policy come Realizzazione di \textbackslash zeta\_\{\textbackslash text\{Sol\}\}(t)}}\label{la-frattura-tech-policy-come-realizzazione-di-zeta_textsolt}

Il Ch.~6 ha modellato il Soliton (Musk) come sorgente esogena di
disordine \(\zeta_{\text{Sol}}(t)\) nel sistema USA. Il decorso fattuale
2025-2026 fornisce verification empirica forte di questa modellazione.

{\def\LTcaptype{none} % do not increment counter
{\small\begin{longtable}[]{@{}
  >{\raggedright\arraybackslash}p{(\linewidth - 4\tabcolsep) * \real{0.3333}}
  >{\raggedright\arraybackslash}p{(\linewidth - 4\tabcolsep) * \real{0.3333}}
  >{\raggedright\arraybackslash}p{(\linewidth - 4\tabcolsep) * \real{0.3333}}@{}}
\toprule\noalign{}
\begin{minipage}[b]{\linewidth}\raggedright
Fase
\end{minipage} & \begin{minipage}[b]{\linewidth}\raggedright
Periodo
\end{minipage} & \begin{minipage}[b]{\linewidth}\raggedright
Effetto su \(\Phi^{\text{USA}}\)
\end{minipage} \\
\midrule\noalign{}
\endhead
\bottomrule\noalign{}
\endlastfoot
Alleanza inaugurale (Zuckerberg, Bezos, Pichai, Musk in posti
privilegiati) & 20 gennaio 2025 & apparente \(\Phi \uparrow\) (mossa di
synchronization visibile) \\
DOGE: 250.000+ posti federali eliminati, target \(2\)T tagli &
gennaio-aprile 2025 & \(\zeta_{\text{Sol}}\) massimo --- distruzione di
sotto-attori istituzionali \\
Magnificent Seven \(-4{,}2\)T (-24\%) tra inauguration e 21 aprile 2025
& Q1 2025 & \(\zeta\) retroattivo: l'alleanza politica produce
decoherence economica \\
Ban su Nvidia H20 chip in China, \(5{,}5\)B charge & Q1 2025 & frizione
interna fra obiettivi politici e settore tecnologico \\
Feud pubblica Musk-Trump $\to$ scuse di Musk & giugno 2025 &
\(\Phi_{\text{cluster Musk-Trump}} \to 0\) \\
Esclusione Musk, Altman, Microsoft dal nuovo tech advisory council;
ingresso Zuckerberg, Ellison, Huang & marzo 2026 & ricomposizione
asimmetrica --- il cluster decoerente si frammenta \\
\end{longtable}}
}

\Hypoth{} \textbf{Tesi formal.} L'equation di Kuramoto modificata del
Ch.~6,
\[\frac{d\theta_j}{dt} = \omega_j + K\Phi \sin(\psi - \theta_j) + \zeta_{\text{Sol}}(t),\]
è ora \emph{empiricamente vincolata}. La sequenza fattuale 2025-2026
mostra che \(\zeta_{\text{Sol}}\) è stato sufficientemente intenso e
persistente da far collassare \(\Phi^{\text{USA}}\) malgrado \(K > K_c\)
formal (il sistema partitico, il sistema federale, il sistema mediatico
esistono ancora). La predizione del Ch.~6.6 --- \emph{correlazione
\(r > 0{,}3\) tra Musk visibility e polarizzazione USA nel 2022-2026}
--- è ora verificationbile sul dataset 2025-2026 con strumenti statistici
standard (vedi Predizione 4 infra).

\Conj{} \textbf{Sotto-tesi non banale.} Il fatto che il \emph{cluster
Musk-Trump} sia esso stesso decoerente entro marzo 2026 è la
realizzazione macro della \textbf{non-trasferibilità del Soliton} (Cap.
6.5): un Soliton non può co-coordinarsi stabilmente con un sistema
egemone in decoherence, because entrambi mancano di \(\Phi\). La
cooperazione produce somma, non interferenza costruttiva.



\section{Tariffe e Decoerenza
Monetaria}\label{tariffe-e-decoherence-monetaria}

\Proved{} \textbf{Dato.} Da gennaio ad aprile 2025 l'aliquota tariffaria
effettiva USA è salita dal \(2{,}5\%\) al \(27\%\) (massimo da oltre un
secolo). Media annua 2025: \(7{,}7\%\), massima dal 1947. Ciò configure
il più grande aumento fiscale come \(\%\) del PIL dal 1993, con costo
medio per famiglia di \(1.500\) nel 2026 (\(1.000\) nel 2025).

\Hypoth{} \textbf{Effetto sistemico DEQ$^2$.} La tariffa generalizzata è
formalmente equivalente a un aumento di \(\sigma\) esogeno per tutti i
partner. Stati con \(A > 0\) assorbono lo shock convessamente; Stati con
\(A < 0\) amplificano. Per gli \textbf{USA stessi}, la tariffa
universale produce auto-shock interno: aumenta \(\sigma\) del proprio
sistema produttivo first di colpire gli interlocutori. È un'azione che
presuppone \(A_{\text{sys}} > 0\) per essere razionale, while la
diagnosi del \S{}7.2 dà \(A_{\text{sys}} < 0\). Si configure therefore come
\textbf{errore strategico structural}, non come scelta antifragile.

\Proved{} \textbf{Reazione del sistema.} La dedollarizzazione, già in
corso da BRICS+, accelera. Non solo governi ma anche istituzioni
finanziarie e investitori privati riducono esposizione al dollaro. La
risposta della Corte Suprema (sentenza 6-3 del 20 febbraio 2026) e la
riformulazione esecutiva via Section 122 (10\% globale per 150 giorni)
testimoniano un sistema che oscilla rapidamente fra vincoli e bypass ---
pattern di \emph{bassa \(\Phi\) con alto throughput}.

\Conj{} \textbf{Predizione DEQ$^2$ (datata).} Quota delle riserve mondiali
in dollari (IMF COFER): \(58{,}4\%\) Q4 2024 \(\to\) scenderà sotto il
\(55\%\) entro Q4 2027 e sotto il \(52\%\) entro Q4 2030. Falsificabile
su dati IMF pubblici trimestrali.



\section{Midterm 2026 come Test di
Falsificazione}\label{midterm-2026-come-test-di-falsification}

\Proved{} \textbf{Dato.} Forecast aggregati ad aprile 2026: probabilità
Dem riprendere la House \(78{,}2\%\); Senato meno favorevole (53-47 GOP,
Dem necessitano \(+4\)). Polymarket assegna \(50{,}5\%\) a \emph{Dem
sweep}, \(36{,}5\%\) a House Dem \(+\) Senato GOP, \(12{,}5\%\) a tenuta
GOP totale.

\Hypoth{} \textbf{Lettura DEQ$^2$: i midterm come misura di \(\Phi\)
residuo.} Se \(\Phi\) fosse ancora alto, il sistema produrrebbe un
riassetto bipartisan --- alternanza fisiologica. Se \(\Phi\) è
collassato, lo \emph{split government} prodotto dai midterm non genera
coordinazione ma \emph{blocco mutuo con escalation}. La predizione
operativa è che lo scenario \emph{Dem House + GOP Senate + Trump
Executive} attivi \(\Omega \uparrow\) (energia dissipata in conflitto
inter-istituzionale) anziché \(\Phi \uparrow\) (coordinazione). I primi
sei mesi del nuovo Congresso (gennaio-luglio 2027) saranno il test
empirico.

\Conj{} \textbf{Predizione testabile.} Numero di vetoes presidenziali
nei primi 12 mesi del Congresso 119$^\circ$ (starting from gennaio 2027):
\(> 15\), contro media storica annua post-1981 di \(\approx 6\). Numero
di sentenze SCOTUS di tipo \emph{emergesncy docket} su questioni
amministrative: \(> 30\) nello stesso periodo, contro media 2015-2020 di
\(\approx 10\).



\section{Chapter Summary}\label{sintesi-del-chapter}

\Proved{} \textbf{In una frase:} the United States del 2026 sono
l'archetipo structural dell'Decoherent Hegemon DEQ$^2$ --- un sistema con
material capacity residua e meccanismi di coupling ancora
formalmente operanti, in cui \(\Phi\) è collassato per effetto
convergente di hyper-extension esecutiva (riduzione di
\(A_{\text{sys}}\)), affective polarization (rumore endogeno
irriducibile), crollo di \(s_{iii}\) (perdita di legitimacy esterna
documentata da Brand Finance e V-Dem) e azione persistente del Soliton
come sorgente esogena \(\zeta(t)\).

{\def\LTcaptype{none} % do not increment counter
{\small\begin{longtable}[]{@{}
  >{\raggedright\arraybackslash}p{(\linewidth - 6\tabcolsep) * \real{0.2500}}
  >{\raggedright\arraybackslash}p{(\linewidth - 6\tabcolsep) * \real{0.2500}}
  >{\raggedright\arraybackslash}p{(\linewidth - 6\tabcolsep) * \real{0.2500}}
  >{\raggedright\arraybackslash}p{(\linewidth - 6\tabcolsep) * \real{0.2500}}@{}}
\toprule\noalign{}
\begin{minipage}[b]{\linewidth}\raggedright
Variabile
\end{minipage} & \begin{minipage}[b]{\linewidth}\raggedright
Valore stimato 2026
\end{minipage} & \begin{minipage}[b]{\linewidth}\raggedright
Direzione
\end{minipage} & \begin{minipage}[b]{\linewidth}\raggedright
Confollowsnza per \(T_s^{(2),\text{sys}}\)
\end{minipage} \\
\midrule\noalign{}
\endhead
\bottomrule\noalign{}
\endlastfoot
\(s_{ii}\) (material capacity) & alta & stabile & mantenuta \\
\(s_{iii}\) (legitimacy) & \(74{,}9\) Brand Finance, \(0{,}57\) V-Dem &
\(\downarrow\downarrow\) & abbatte numeratore \\
\(\Phi\) (phase coherence) & \(\to 0\) & crollo structural & annulla
numeratore \\
\(A_{\text{sys}}\) (antifragility) & \(< 0\) & concavo & amplifica
\(\sigma\) \\
\(\Omega\) (Neijuan) & \(\uparrow\) & conflitto interno & gonfia
denominatore \\
\(\varepsilon_{\text{dis}}\) (disimpegno) & \(\uparrow\) (40\%+ vede
l'altro come minaccia) & \(\uparrow\) & gonfia denominatore \\
\(\zeta_{\text{Sol}}\) (sorgente esogena) & persistente, anche post-feud
& \(\uparrow\) & forza decoherence \\
\end{longtable}}
}

\Hypoth{} \textbf{Tesi conclusiva del chapter.}
\(T_s^{(2),\text{sys}\,,\text{USA}} \to 0\) entro la finestra
\textbf{2027-2029}, con probabilità superiore a \(0{,}65\) conditional
agli indicatori 2026. La decoherence è asimmetrica with respect tol collapse:
gli USA \emph{non} perdono la material capacity (\(s_{ii}\) resta
alto), ma perdono la \textbf{funzione egemonica} --- la possibilità di
proiettare un order parameter globale. Restano una \emph{potenza
materiale senza fase}, equivalente sistemico di una stella che continua
a brillare while il suo nucleo si è già spento.



\section{Tre Predizioni Datate del Ch.~7 (per Cap.
11)}\label{tre-predizioni-datate-del-cap.-7-per-cap.-11}

\begin{enumerate}
\def\labelenumi{\arabic{enumi}.}
\tightlist
\item
  \textbf{V-Dem USA \(< 0{,}55\) entro report 2027} (rilascio aprile
  2027). Falsificabile su \url{https://v-dem.net}.
\item
  \textbf{Quota dollaro nelle riserve IMF COFER \(< 55\%\) entro Q4 2027
  e \(< 52\%\) entro Q4 2030.} Falsificabile su dati IMF trimestrali.
\item
  \textbf{Vetoes presidenziali \(> 15\) nel primo anno del Congresso
  119$^\circ$ (gennaio-luglio 2027); emergesncy docket SCOTUS \(> 30\) stesso
  periodo.} Falsificabile su Federal Register e statistics SCOTUSblog.
\item
  \textbf{Correlazione \emph{Musk visibility}--polarizzazione USA
  \(r > 0{,}3\) nel periodo 2022-2026} (estensione predizione Ch.~6.6,
  calibrabile ora su dati post-DOGE).
\end{enumerate}



\section{Closing Epistemic Notes
Chapter}\label{note-epistemiche-di-chiusura-chapter}

{\def\LTcaptype{none} % do not increment counter
{\small\begin{longtable}[]{@{}
  >{\raggedright\arraybackslash}p{(\linewidth - 4\tabcolsep) * \real{0.3333}}
  >{\raggedright\arraybackslash}p{(\linewidth - 4\tabcolsep) * \real{0.3333}}
  >{\raggedright\arraybackslash}p{(\linewidth - 4\tabcolsep) * \real{0.3333}}@{}}
\toprule\noalign{}
\begin{minipage}[b]{\linewidth}\raggedright
\#
\end{minipage} & \begin{minipage}[b]{\linewidth}\raggedright
Claim
\end{minipage} & \begin{minipage}[b]{\linewidth}\raggedright
Status
\end{minipage} \\
\midrule\noalign{}
\endhead
\bottomrule\noalign{}
\endlastfoot
1 & Definizione formal Decoherent Hegemon & \Hypoth{} categoria DEQ$^2$
nuova \\
2 & Crollo V-Dem 2024$\to$2026 (\(0{,}73 \to 0{,}57\)) & \Proved{} dato
pubblicato \\
3 & Brand Finance \(-4{,}6\) pt (calo più ripido di tutti) & \Proved{}
dato pubblicato 2026 \\
4 & Iper-estensione \(\Rightarrow A_{\text{sys}} < 0\) & \Hypoth{}
inferenza structural \\
5 & Solitonizzazione discorsiva degli USA \((7,6,8) \to (8,9,2)\) &
\Conj{} mappatura interpretativa, falsificabile su corpus
presidenziali \\
6 & \(\zeta_{\text{Sol}}\) verified fattualmente (DOGE, feud, cluster
decoerente) & \Proved{} evidenza convergente \\
7 & Tariffa universale come errore strategico per \(A < 0\) & \Hypoth{}
confollowsnza followsta \\
8 & Predizione \(T_s^{(2),\text{sys}} \to 0\) in 2027-2029,
\(P > 0{,}65\) & \Conj{} conjecture datata \\
9 & Quota dollaro \(< 55\%\) entro Q4 2027 & \Conj{} predizione
testabile \\
10 & Vetoes \(> 15\) + emergesncy docket \(> 30\) in un anno & \Conj{}
predizione testabile \\
\end{longtable}}
}



\section{Fonti dei dati 2025-2026
utilizzate}\label{fonti-dei-dati-2025-2026-utilizzate}

\begin{itemize}
\tightlist
\item
  V-Dem Democracy Report 2026 --- declassamento USA a \emph{electoral
  democracy}, score \(0{,}57\)
\item
  Pew Research, dicembre 2025 --- trust nel governo federale Dem \(9\%\)
  / Rep \(26\%\)
\item
  Brand Finance Global Soft Power Index 2026 --- calo \(-4{,}6\) punti,
  dettaglio per asse
\item
  Ballotpedia / Federal Register --- count Executive Orders Trump 2.0 al
  31 marzo 2026
\item
  Tax Foundation --- aliquota effettiva tariffaria \(27\%\) aprile 2025,
  media \(7{,}7\%\) 2025
\item
  Supreme Court of the United States --- \emph{Trump v. CASA}
  27/06/2025, sentenza tariffe IEEPA 20/02/2026
\item
  Race to the WH / 270toWin / Polymarket --- forecast midterm 2026
\end{itemize}


